% =====================================================================
% Preâmbulo: pacotes e formatação do relatório
% =====================================================================
\usepackage[utf8]{inputenc}
\usepackage[T1]{fontenc}
\usepackage[brazil]{babel}
\usepackage{mathptmx}
\usepackage{amsmath,amssymb}
\usepackage[a4paper,margin=2.5cm,headheight=14.5pt]{geometry}
\usepackage{setspace}
\usepackage{indentfirst}
\usepackage{microtype}
\usepackage{graphicx}
\usepackage{booktabs}
\usepackage{array}
\usepackage{caption}
\usepackage{float}
\usepackage{enumitem}
\usepackage{tikz}
\usepackage{newfloat}
\usepackage{fancyvrb}
\usepackage{fvextra}
\usepackage{fancyhdr}
\usepackage[hyphens]{url}
\usepackage[round,semicolon,authoryear]{natbib}
\usepackage[hidelinks]{hyperref}

\onehalfspacing
\frenchspacing
\setlength{\parindent}{1.25cm}
\setlength{\parskip}{0pt}
\raggedbottom
\setlist{itemsep=0.25em,topsep=0.4em}
\captionsetup{font=small,labelfont=bf,skip=6pt}
\DeclareFloatingEnvironment[fileext=loa,listname={Lista de algoritmos},name=Algoritmo,placement=htbp]{algoritmo}
\graphicspath{{figuras/}}
\urlstyle{rm}

\pagestyle{fancy}
\fancyhf{}
\fancyhead[R]{\thepage}
\renewcommand{\headrulewidth}{0pt}
\fancypagestyle{plain}{%
  \fancyhf{}%
  \fancyhead[R]{\thepage}%
  \renewcommand{\headrulewidth}{0pt}%
}
\numberwithin{equation}{section}

\makeatletter
\renewcommand{\@biblabel}[1]{}
\makeatother
\setlength{\bibhang}{1.25cm}

\title{do Cálculo Numérico e sua Relação com a Ciência da Computação: Estudo e Implementação dos Métodos de Runge-Kutta na Simulação de Circuitos RLC}

\date{2026}
